\section{The analytic series and the pair comparison}\label{sec:pair}

Fix an odd prime \(p\).  We retain the imaginary quadratic field \(K\),
the auxiliary field \(L\), and the tuple of tame directions constructed
in Sections~\ref{sec:classical} and~\ref{sec:vertical}.  Thus \(p=w\bar w\)
in \(K\), every rational prime below the fixed finite support for
\(E/K\) splits in \(K\), and
\[
 R=\Z_p[[t,u_1,\ldots,u_d]]
\]
has one genuine anticyclotomic variable \(t\).  Its group is the maximal
free pro-\(p\) quotient \(\Gamma\simeq\Z_p\) of the ring-class tower at
\(p\).  The variables \(\mathbf u=(u_1,\ldots,u_d)\) arise from the
moving tame conductors.  Write \(L_w\in R\) for the determinant of the
strict-at-\(w\), full-at-\(\bar w\) complex, with full conditions at the
non-\(p\) places of the support.  It is defined up to an integral unit.
The tuple satisfies
\[
 L_w\not\equiv0\pmod{(p,t)}
\]
for either choice of \(w\), and likewise for the companion curve, by
\eqref{eq:source-1}.

Our aim here is to construct a bounded analytic series \(B_w\), with its
exact central value, and to prove that \(B_w/L_w\) is integral and has
no factor \(p\).  We will also compare the two choices of \(p\)-adic
place at finite characters.  These statements retain possible
horizontal divisors of the quotient; their removal is the task of
Section~\ref{sec:unit}.

Choose a modular parametrization \(\phi:X_0(N)\longrightarrow E\),
based at the cusp \(\infty\), and write
\[
 \phi^*\omega_E=c_E f\,dq_{\rm Tate}/q_{\rm Tate}.
\]
The Manin factor \(c_E\) will remain in every exact formula.  Let
\(S^\circ\) be the fixed set of rational primes away from \(p\) in the
\(E/K\) comparison.  These primes split in \(K\).  Enlarging the fixed
support for the companion is allowed; it does not change the meaning
of \(S^\circ\) for \(E\).

The arguments below include CM curves.  The characteristic-zero
representations \(V|_{G_K}\) and the analogous representation for the
companion are absolutely simple.  For a non-CM curve this follows from
the open image theorem \cite{Serre}.  For a CM curve, \(K\) and \(L\)
were chosen to avoid the CM field.  The two inducing characters remain
distinct on the compositum, as is seen from their CM types, or from
the distinct weights of their \(p\)-adic avatars.  Restriction therefore
preserves absolute simplicity in this case also.

\subsection{Integral primitives on the ordinary locus}

For every rational prime \(q\), put
\[
 P_q(Z)=1-\frac{a_q}{q}Z+\frac{\epsilon_q}{q}Z^2,\qquad
 \epsilon_q=
 \begin{cases}
 1,&q\nmid N,\\
 0,&q\mid N.
 \end{cases}
\]
At a split prime away from \(p\), choose an oriented cyclic level,
raised as necessary to define the quotient operator \(V_q\).
Normalize \(V_q\) by the substitution \(q_{\rm Tate}\mapsto
q_{\rm Tate}^q\).  At \(p\) use only the canonical multiplicative
cyclic level.  On a CM disk at \(w\), let
\[
 G(X)=c_E^{-1}\log_{\omega_E}\phi(X).
\]
Let \(d_{\rm mod}\) be ordinary modular differentiation: on Tate
charts it is \(q_{\rm Tate}\partial_{q_{\rm Tate}}\), and on the
Igusa tower it is dual to the Kodaira--Spencer differential in the
square of the trivialized invariant differential.
The following construction adapts the ordinary-locus arguments of
\cite[Section~8, \emph{Integral depleted primitive},
\emph{Support of the disk measure}, and
\emph{The exact logarithmic primitive}]{B}.
We give the argument at odd \(p\), including the case \(p\mid N\).
Here \emph{ordinary} refers to the elliptic curves parametrized by
the modular source.  It imposes no ordinary-reduction hypothesis
on the target \(E\): the depletion and inverse-derivative limit
use the coefficients \(a_n\), without choosing a unit root.

\begin{proposition}[Integral logarithmic primitive]\label{pair:primitive}
On the canonical ordinary locus define
\[
 f^{[p]}=(1-a_pV_p+p\epsilon_pV_p^2)f,\qquad
 F_0=\lim_{j\to\infty}d_{\rm mod}^{(p-1)p^j-1}f^{[p]},\qquad
 F=\prod_{q\in S^\circ}P_q(V_q)F_0.
\]
The limits are integral weight-zero functions.  They satisfy
\[
 d_{\rm mod}F_0=f^{[p]},\qquad
 F=\sum_{\substack{n\ge1\\(n,p\prod_{q\in S^\circ}q)=1}}
          \frac{a_n}{n}q_{\rm Tate}^n
\]
at a Tate cusp with all cyclic groups of multiplicative type.
If \(z\) is a Serre--Tate coordinate with \(z=1\) at the canonical
lift, the expansion \(F(z)\) is the transform of a bounded integral
measure \(\mu\) on \(\Z_p\), supported on \(\Z_p^\times\).
On every interior CM disk one has the exact identity
\begin{equation}
 F_0=P_p(V_p)G.
 \tag{9}\label{eq:source-9}
\end{equation}
\end{proposition}

\begin{proof}
Work first on a fine full prime-to-\(p\) level over
\(\mathcal V=W(\overline{\F}_p)\).  The ordinary elliptic moduli and
Serre--Tate theory \cite{DeligneRapoport,KatzMazur} provide the
canonical connected subgroup and the
Igusa tower trivializing the multiplicative formal group.  The dual of
the Kodaira--Spencer differential, expressed using the square of the
trivialized invariant differential, defines \(d_{\rm mod}\); on Tate
charts it is \(q_{\rm Tate}\partial_{q_{\rm Tate}}\).
For quotient maps we use the target differential whose pullback has
the corresponding degree factor.  These are precisely the
ordinary-moduli conventions used in the cited lemmas.

The initial algebraic pullbacks have bounded denominators on the
quasi-compact ordinary formal region, including its Tate charts.
Integrality can then be tested by successive division by the
uniformizer and expansion at cusps.  The special-fiber determinant
components of the fine tame curve are irreducible, and every
component of a finite trivialization level surjects onto its ordinary
component.  Thus expansions at all trivialization lifts of a
multiplicative cusp detect zero reduction on every component.
Modulo any fixed precision, a section on the completed tower uses a
finite level, so the same test applies there.

The \(p\)-depletion has no terms whose index is divisible by \(p\).
For \(n\in\Z_p^\times\), the numbers
\(n^{(p-1)p^j-1}\) tend uniformly to \(n^{-1}\).
The weights \(2(p-1)p^j\) tend to zero on
\(\Z_p^\times\).  Consequently the differentiated expansions converge
integrally on every trivialization.  The expansion test just proved
gives convergence on the tower.  The limit is invariant under the
full \(\Z_p^\times\)-action on differential trivializations, so
torsor descent gives \(F_0\) on the ordinary formal curve at the
fixed fine tame level.  Its reduction modulo \(p\) is consequently
a regular function on that curve's ordinary special fiber.
The Hecke recursions give the displayed
fully depleted expansion of \(F\).  This construction uses the
canonical multiplicative \(p\)-structure, so it also applies when
\(p\mid N\).

Mahler duality identifies the integral expansion in \(z-1\) with a
bounded measure.  The trace of \(F\) along the canonical quotient is
zero: on the parameter it is the finite flat map \(z\mapsto z^p\),
reducing to Frobenius, and at a Tate cusp the trace kills every
prime-to-\(p\) exponent.  Expansion detection therefore gives
\[
 \sum_{\zeta^p=1}F(\zeta z)=0.
\]
In the measure transform this sum multiplies the restriction to
\(p\Z_p\) by \(p\) and kills its complement.  Uniqueness of transforms
and torsion-freeness show that this restriction is zero.

It remains to identify the constant in the antiderivative.
The map to \(E\) extends by the Neron mapping property on the smooth
ordinary charts carrying the canonical \(p\)-subgroup.  Such a chart
has the same deformation germ as its tame chart: a subgroup lifting
the connected special subgroup has zero image in the etale quotient
and must be the canonical multiplicative subgroup.  In a Neron
residue tube the formal group logarithm, translated by the value at
the center, is analytic.  Multiplication into a smaller formal
neighborhood and division by that integer gives the same analytic
logarithm throughout the interior disk.  This proves analyticity of
\(G\), including when the elliptic target has bad reduction.

The differentials of \(F_0\) and \(P_p(V_p)G\) agree.  Indeed
differentiating \(V_pG\) introduces a factor \(p\), and
differentiating \(V_p^2G\) introduces \(p^2\).  Their difference is
therefore constant on the disk.  The noncanonical quotient sources
of a target \(B\) give, by the \(T_p\) or \(U_p\) correspondence,
\[
 \sum_{\zeta^p=1}G(\zeta z)
       =a_pG(B)-\epsilon_pG(V_pB).
\]
Cuspidal differences are torsion \cite{Manin,ManinDrinfeld} and
disappear under the logarithm.
When \(p\mid N\), the noncanonical quotients retain the required
canonical level, so the same identity uses \(U_p\).
It follows that \(P_p(V_p)G\), like \(F_0\), has trace zero.
The \(p\) sources lie in the same residue disk.  The trace of their
constant difference is consequently \(p\) times that difference,
which is zero in characteristic zero.  This proves
\eqref{eq:source-9}.  Equivalently the same normalization is the
elliptic weight-two case of \cite[Section~5, \emph{Raw toric periods
and the inverse Katz derivative}]{C}; the preceding argument specifies its
extension to a potentially bad target and to the canonical
\(p\)-structures used here.
\end{proof}

\subsection{CM orbits and exact interpolation}

At a finite stage let \(s\) be the product of the moving tame
conductors; the case \(s=1\) is included.  The canonical centers have
endomorphism order \(\mathcal O_s\).  At a primitive root
\(z=\zeta_{p^m}^{\,j}\), \(j\in(\Z/p^m)^\times\), endomorphism
lifting gives order \(\mathcal O_{p^ms}\).  The point, with its cyclic
level, is defined over the ring-class field \(H_{p^ms}\):
the completed order units preserve the canonical \(p\)-level and
the oriented prime-to-\(p\) level.  Here \(\mathcal O_K^\times=\{\pm1\}\).

The kernel of the order-class map to conductor \(s\) acts simply
transitively on the primitive roots in each disk.  In Serre--Tate
coordinates its action is the ratio of the connected and etale local
unit actions.  The order-class exact sequence thus gives a map
\[
 \eta_\Gamma:\Z_p^\times\longrightarrow\Gamma
\]
independent of \(s\), with kernel \(\mu_{p-1}\) and open image.
We use the actual Artin labels for this map.  Inverting all reciprocity
conventions simultaneously makes no change to the assertions below.
These coordinate facts are the odd-\(p\) version of
\cite[Section~8, \emph{CM disk coordinates}]{B}:
the endomorphism-lifting condition is equality of the two local
actions modulo \(p^m\), and the corresponding class kernel is
\((\Z/p^m)^\times\).

Choose compatible ideals \(\mathfrak a\), prime to conductors and
levels, representing \(\Pic(\mathcal O_s)\).  Let
\(t_{\mathfrak a}\in\Gamma\) be their compatible Artin values and
\(\delta_{\mathfrak a}\in\Z_p^\times\) the exponents that transport
the Serre--Tate frames from the base disk.  Let
\(\psi_s\) be the tuple character in the finite tame group rings.
For \(y\in\Gamma\) write \([y]\) for its group-like power of \(1+t\).
If \(\mu_{\mathfrak a}\) is the measure in
Proposition~\ref{pair:primitive} on the disk indexed by \(\mathfrak a\),
define
\[
 b_s(t,\mathbf u)=
 \sum_{\mathfrak a\in\Pic(\mathcal O_s)}
 [\,t_{\mathfrak a}-\eta_\Gamma(\delta_{\mathfrak a})\,]\,
 \psi_s(\mathfrak a)
 \int_{\Z_p^\times}[-\eta_\Gamma(x)]\,d\mu_{\mathfrak a}(x).
\]
This is a bounded integral series with finite group-ring tame
coefficients.  Form it with two opposite non-\(p\) orientations,
using the same canonical \(p\)-level and underlying disks.
Passing to the bounded coefficient limits of Section~\ref{sec:models}
gives \(b_w^+\), \(b_w^-\), and
\(B_w=b_w^+b_w^-\).  The coefficient ring is initially
\(\mathcal V^*[[t,\mathbf u]]\), where
\(\mathcal V^*=W(\overline{\F}_p^{\,*})\) is the unramified
coefficient enlargement defined there.

\begin{proposition}[Finite characters and the center]\label{pair:interpolation}
Let \(\theta\) be a sufficiently high finite character of \(\Gamma\),
and let \(\theta_l=\theta\circ\eta_\Gamma\) have primitive conductor \(p^m\).
Put \(\chi=\theta\psi_s\), and let \(X\) be the point with parameter
\(\zeta_{p^m}\) in the base disk at the original primitive level.
Then
\begin{equation}
 \left(\sum_{j\bmod p^m}\theta_l(j)\zeta_{p^m}^{\,j}\right)
     b_s(\theta,\mathbf u)
 =
 \left(\sum_{\rho\in\Gal(H_{p^ms}/K)}
        \chi(\rho)G(X^\rho)\right)
        \prod_{q\in S^\circ}P_q(\chi(\sigma_q)).
 \tag{10}\label{eq:source-10}
\end{equation}
The character is extended by zero on nonunits in the Gauss sum, and
\(\sigma_q\) is inverse translation by the oriented \(V_q\)-action.
Let \(b_w^0\) denote the first limiting series formed with \(F_0\),
without the non-\(p\) depletions.  Then, over the integral coefficient
enlargement,
\begin{equation}
 \begin{split}
 B_w&=\text{\rm unit}\cdot(b_w^0)^2D,\\
 D&=\prod_{q\in S^\circ}
 P_q(\chi(\Frob_{v_q}))
 P_q(\chi(\Frob_{v_q})^{-1}),
 \end{split}
 \tag{11}\label{eq:source-11}
\end{equation}
where in this identity \(\chi\) is the universal power-series
character and either place \(v_q\mid q\) may be chosen.
For conductor \(s=1\), the central value is
\begin{equation}
 B_w^{s=1}(0)=
 \pm c_E^{-2}\log_{\omega_E}(P_K)^2P_p(1)^2
       \prod_{q\in S^\circ}P_q(1)^2,
 \tag{12}\label{eq:source-12}
\end{equation}
where \(P_K\) is the single Hilbert-class trace of the conductor-one
Heegner point under \(\phi\).
\end{proposition}

\begin{proof}
Finite Fourier summation on \((\Z/p^m)^\times\) multiplies the inverse
moment of a measure by the displayed Gauss sum.  Transporting the
frame in disk \(\mathfrak a\) changes the root exponent by
\(\delta_{\mathfrak a}\).  The correction
\([t_{\mathfrak a}-\eta_\Gamma(\delta_{\mathfrak a})]\) therefore makes its
weight exactly the Artin character \(\chi\).
The class-kernel description identifies the resulting sum over roots
and disk centers with the sum over \(\Gal(H_{p^ms}/K)\).

Apply \eqref{eq:source-9}.  For sufficiently large \(m\), the \(V_p\)
and \(V_p^2\) terms factor through lower-conductor orbits, so their
sums against the primitive character vanish.  Each remaining
quotient \(V_q\) translates the orbit by \(\sigma_q\), with its level
choices retained.  This proves \eqref{eq:source-10}, with no averaging
factor.

Atkin--Lehner at the non-\(p\) part of the primitive level reverses
the non-\(p\) orientations.  On the newform quotient it acts by a
sign; the change of base cusp is torsion.  The raw logarithmic sum
therefore changes only by that sign and a translation character.
The same assertion holds when a subset of the non-\(p\) orientations
is reversed.  The translating ideals are prime to \(p\), so these
translations can be chosen compatibly in the \(p\)-tower.  Applying
\eqref{eq:source-10} at all high characters, followed by uniqueness
for bounded series, gives \eqref{eq:source-11}.  At the moving stages
the group-like translations likewise have coefficientwise limits.

At the trivial character, integration is evaluation at the disk
center.  Formula \eqref{eq:source-9} now retains the \(p\)-Euler
factor, because \(V_p\) acts by translation on the conductor-one
centers.  The sum is taken once over the class group.  The two
orientation sums are translates of the same Hilbert trace up to
sign.  Recalling \(G=c_E^{-1}\log_{\omega_E}\phi\) gives
\eqref{eq:source-12}.
\end{proof}

\begin{lemma}[Removal of tame conductor variables]\label{pair:trace}
Choose each new base disk by a descending isogeny from the preceding
conductor.  Setting a new tame variable equal to zero multiplies
each single analytic series by
\begin{equation}
 a_r-\chi_{\rm old}(\Frob_{v_r})
           -\chi_{\rm old}(\Frob_{v_r})^{-1}.
 \tag{13}\label{eq:source-13}
\end{equation}
The notation denotes the limit of the stage coefficients and
character values.  The determinant \(L_w\) is multiplied by the
square of this factor up to an integral unit.  These assertions hold
successively for any subset of the new variables, including removal
of all tame variables.
\end{lemma}

\begin{proof}
The good Hecke correspondence at \(r\), after subtracting its two
horizontal split translations, is the trace from conductor \(rs\)
to conductor \(s\).  Its action on the character sum in
\eqref{eq:source-10} is \eqref{eq:source-13}.
Transport the frames by polarization.  Their exponents are
\(r_i^{\pm1}\), which tend to \(1\) modulo every fixed \(p^m\).
Thus the same identity holds at each fixed sufficiently high
\(\theta\), simultaneously in the group rings of the tame variables
that remain.  The expressions defining \(b_s\) are bounded as group
polynomials.  We can consequently pass to limits in those group
rings; bounded character projectors are not required.

For the arithmetic determinant, unramified inflation restores the
full pair over each removed prime.  The two singular determinants
are the unramified Tate Euler polynomials at that pair.
Since \(r_i\to1\), their product is the square of
\eqref{eq:source-13} up to an integral unit.  This is an identity of
the marked integral determinant triangles, not a comparison after
rationalization.  Iterating with the old tuple fixed proves all the
stated specializations.  The factor need not be a unit at the
central specialization.
\end{proof}

\begin{lemma}[Descent of the analytic ideals]\label{pair:descent}
The principal ideals generated by \(b_w^+\), \(b_w^-\), \(b_w^0\),
and \(B_w\) descend to \(R\).  Each actual series differs from a
generator over \(R\) by an integral unit.
\end{lemma}

\begin{proof}
Extend arithmetic unramified Frobenius \(\varphi\) at \(w\) to fix
the cyclotomic \(p\)-power roots.  At every finite stage,
\eqref{eq:source-10} identifies its action on a raw CM sum with
translation by a fixed local Artin element: the ring-class extension
is abelian over \(K\), so this action commutes with all labels.
The Gauss coordinates are fixed.  Passing to limits and testing at
high finite characters gives, for each series in question,
\[
 \varphi(b)=\kappa b,
\]
where \(\kappa\) is group-like of augmentation \(1\).
On \(\mathcal V^*\), \(\varphi\) is Witt Frobenius.
There is a unit \(e\), with constant coefficient \(1\), satisfying
\(\varphi(e)=\kappa e\).  Solve recursively in total degree: at each
step the remaining equation is \(\varphi(a)-a=b_0\).
It is soluble modulo \(p\) because the residue field is algebraically
closed, and successive lifting solves it in \(\mathcal V^*\).
Then \(b/e\) has coefficients fixed by Witt Frobenius, hence in
\(\Z_p\).  Products give the assertion for \(B_w\).
We continue to use the actual analytic series for value formulas
and the descended generators for height-one ideal comparisons.
\end{proof}

\subsection{Primitivity of the analytic series}

The preceding construction controls the period and Euler factors.
The next argument provides the separate vertical input: no fixed
power of \(p\) is lost in the analytic series.  We prove the required
tame-stalk independence by ordinary deformation theory and monodromy.

\begin{proposition}[Vertical analytic primitivity]\label{pair:vertical-primitivity}
For each of the tuples under consideration and either \(p\)-adic
place,
\begin{equation}
 p\nmid B_w.
 \tag{14}\label{eq:source-14}
\end{equation}
\end{proposition}

\begin{proof}
We first establish the independence statement in the elliptic case.
Let \(k=\overline{\F}_p\), let \(x\) be an ordinary maximal-order
CM point with endomorphism ring \(\mathcal O_K\), and let \(\mathscr R\)
be its algebraic local ring on the prime-to-\(p\) modular tower.
Thus a finite set of elements of \(\mathscr R\) comes from a common
finite tame level; we do not replace \(\mathscr R\) by its completion.
For \(a_1,\ldots,a_h\in\Z_p^\times\), consider
\[
 \Phi:\mathscr R^{\otimes_k h}\longrightarrow k[[z-1]],\qquad
 f_1\otimes\cdots\otimes f_h\longmapsto
       \prod_{i=1}^h f_i(z^{a_i}).
\]
We will prove that \(\Phi\) is injective when the \(a_i\) have
pairwise inequivalent ratios under
\[
 \{\alpha/\bar\alpha:
        \alpha\in K^\times,\ \alpha\text{ a unit at both primes over }p\}.
\]

Suppose the prime kernel of \(\Phi\) is nonzero and choose \(h\)
minimal.  Let \(X\) be its irreducible relation locus in the product
of the tame towers.  Its projections to every proper product of
factors are dominant by minimality.  The common CM stabilizer
preserves the kernel: its scalar Serre--Tate substitutions commute
with all the \(a_i\).  Weak approximation in \(K\) makes the ratios
of its two local unit actions dense in \(\Z_p^\times\).  The local
rigidity and normalization argument
of \cite[Proposition~3.11 and Corollary~3.13]{HidaMu} therefore gives
a finite tame-level model and a normalized branch over
\((x,\ldots,x)\) whose completion is a formal subtorus.  This branch
contains \((z^{a_1},\ldots,z^{a_h})\).  Write \(Y\) for the
normalization of the irreducible model.  Dominance onto \(h-1\)
factors and properness give \(\dim Y=h-1\).

Here is the global constraint on that formal subtorus.  Pull back the
product of the universal elliptic schemes to \(Y\).  The global
canonical-coordinate map of \cite[Sections~2 and~5]{ChaiOrd} has,
on a normal base, a cotorsion-free smooth kernel.  Its annihilator
\(\mathscr N\) is a smooth cocharacter subsheaf.  Its rank can be
evaluated on the chosen formal branch, where it is \(h-1\).
The proof of \cite[Proposition~5.3]{ChaiOrd}, applied to the
normalization and this branch, consequently gives a dense smooth
open \(U\subset Y\) on which the same rank is \(h-1\).
The elliptic-factor idempotents of the product abelian scheme place
\(\mathscr N\) inside the direct sum of its \(h\) diagonal
Serre--Tate cocharacter lines.  Each coordinate projection of
\(\mathscr N\) is nonzero.  Otherwise that elliptic deformation
would be formally constant on \(U\), contrary to the dominance of
the corresponding projection of \(X\).  This construction uses the
canonical-coordinate relation sheaf on a normal base; the ranks of
two-coordinate projections alone do not determine its rank.

Descend \(Y\), \(U\), the elliptic schemes, and their
canonical-coordinate sheaves to a finite field.  Suppose first that
the \(h\) generic elliptic factors are pairwise
nonisogenous over an algebraic closure.  Choose an odd prime
\(\ell\ne p\).  After a finite etale cover of a smaller open
\(U\), all generic endomorphisms are defined and the joint
arithmetic algebraic \(\ell\)-adic monodromy group is connected.
The semisimplicity and isogeny theorems apply to its finitely
generated function field \cite[Theorems~1.2 and~1.4]{Zarhin}.
Each nonconstant elliptic factor
has derived monodromy \(\mathrm{SL}_2\).  The Lie algebra of the
derived group is a subdirect product of these simple nonabelian
Lie algebras.  The Lie-algebra Goursat argument decomposes it into
diagonal blocks: the kernel of a projection is an ideal, and a
simple factor is either independent or identified with one of the
preceding simple factors.  A block meeting two factors identifies
their standard two-dimensional representations on the derived
group.  The connected center acts on these representations by
scalar characters whose squares agree, because their determinants
are equal.  The two characters therefore agree.  The intertwiner
is equivariant for the whole connected group, and the isogeny
theorem over finitely generated fields of characteristic \(p\)
makes the two generic elliptic curves isogenous.  Thus no such block
is possible, and the derived monodromy group is
\(\mathrm{SL}_2^h\).  Passing to the connected group has also
removed possible finite quadratic-twist correlations.  These are
the endomorphism and monodromy inputs used in
\cite[proof of Corollary~3.16]{HidaMu}.

Chebotarev gives a closed point \(u\in U\) at which the elliptic
Frobenius traces
\(c_1,\ldots,c_h\) have pairwise distinct squares: the equations
\(c_i^2=c_j^2\) define proper closed subsets of the monodromy
group, whose \(h\) derived factors vary independently.
Let \(\alpha_i\) be the \(p\)-adic unit root at \(u\), and put
\(q_u=\#k(u)\).  If \(\alpha_i^2=\alpha_j^2\), then
\[
 c_i=\alpha_i+q_u/\alpha_i
       =\pm(\alpha_j+q_u/\alpha_j)=\pm c_j,
\]
a contradiction.  Arithmetic Galois descent Frobenius acts on the
\(i\)-th Serre--Tate character line by \(\alpha_i^2\), since that
line is the tensor square of the etale Tate line; it acts on the
cocharacter line by \(\alpha_i^{-2}\).  These are distinct
eigenvalues.  The stable subspace \(\mathscr N_u\otimes\Q_p\)
must therefore be a sum of coordinate axes: each eigenprojection
is a polynomial in Frobenius.  Its nonzero projection to every
coordinate forces it to contain every axis, contradicting its
rank \(h-1\).

There is consequently a generically isogenous pair.  Extend its
isogeny over a normal finite cover of \(Y\); homomorphisms between
abelian schemes extend over a normal base.  At a point above the
chosen CM branch its specialization is a nonzero
\(\alpha\in\operatorname{End}^0(x)=K\).  Serre--Tate functoriality
and the formal curve in that branch give
\(a_i/a_j=\alpha/\bar\alpha\), up to simultaneous inversion.
One may see this before inverting \(p\) by clearing denominators
in the two local actions of \(\alpha\) and comparing exponents in
\(k[[z-1]]\); the map \(b\mapsto z^b\) from \(\Z_p\) is
injective.  Since \(a_i/a_j\) is a unit, the two valuations of
\(\alpha\) above \(p\) agree.  Rational scaling makes both zero
without changing their ratio.  This contradicts the stipulated
inequivalence and proves injectivity of \(\Phi\).  In particular,
a sum of functions from distinct slots is constant only if every
slot function is constant.

We now work at \(s=1\), with one of the two non-\(p\) orientations.
The functions here are regular weight-zero ordinary functions at a
maximal-order CM point.  The weight-zero torsor descent in
Proposition~\ref{pair:primitive} places their reductions on the tame
ordinary special fiber, not merely on a finite Igusa cover.
The canonical multiplicative \(p\)-structure has the same deformation
germ as this tame curve.  Their germs therefore belong to the
algebraic prime-to-\(p\) stalks just considered.

Suppose \(b_1\) is zero modulo \(p\).  Restrict its measure to the
open subgroup \(\eta_\Gamma(\Z_p^\times)\) of \(\Gamma\).
For each contributing class choose
\(\lambda_{\mathfrak a}\in\Z_p^\times\) with
\(\eta_\Gamma(\lambda_{\mathfrak a})=t_{\mathfrak a}\).
Pushing the measure through \(\eta_\Gamma\) folds it over its kernel
\(\mu_{p-1}\).  On transforms the assumed vanishing becomes
\[
 \sum_{\mathfrak a}\sum_{\zeta\in\mu_{p-1}}
 F_{\mathfrak a}
     \bigl(z^{\zeta\delta_{\mathfrak a}/\lambda_{\mathfrak a}}\bigr)
       =0\pmod p.
\]
The class sum includes exactly those classes that contribute to
the indicated open subgroup.  Pull back by the ideal isogenies that
transport the frames to a common prime-to-\(p\) level.  This absorbs
\(\delta_{\mathfrak a}\) into the pulled-back function, so the
exponent whose slot is to be tested is
\(\zeta/\lambda_{\mathfrak a}\).  Combine equivalent slots before
applying the injectivity of \(\Phi\).
An equivalence \(\alpha/\bar\alpha\) is realized by the germ of a CM
quasi-isogeny.  Since this ratio is a \(p\)-unit, rational scaling
makes the isogeny prime to \(p\).

We describe the slot equivalent to \(1\).  Put \(h=h_K\) and
\(\mathfrak a^h=(x)\).  The class-sequence description gives
\[
 \lambda_{\mathfrak a}^{\,h}
       =(x/\bar x)\xi
\]
up to simultaneous inversion, where \(\xi\) is a torsion unit.
If \(\zeta/\lambda_{\mathfrak a}=\alpha/\bar\alpha\), raise to a
power killing \(\xi\) and \(\zeta\), and compare ideal valuations.
The group of fractional ideals is torsion-free; cancellation of
that power and of \(h\) shows that a principal scaling of
\(\mathfrak a\) is conjugation invariant.  This is cancellation in
the fractional ideal group, not division by \(h\) in \(\Z_p\).
Such a class is represented by a product of ramified prime ideals,
each used once or omitted.

Conversely, choose those ramified products as representatives.
Their squares are rational principal ideals, so their images
\(t_{\mathfrak a}\) in the torsion-free group \(\Gamma\) are zero;
we may take \(\lambda_{\mathfrak a}=1\).
Only \(\zeta=\pm1\) can belong to the slot of \(1\), since a torsion
element of \(K^\times\) is then a sign.
The ramified products represent these classes twice, paired by
complementary subsets.  To check the multiplicity, if a ramified
product is \((\alpha)\), then \(\alpha/\bar\alpha=\pm1\).
The plus sign makes \(\alpha\) rational and forces the subset to be
empty; since the discriminant is odd and fundamental, the minus
sign makes it a rational multiple of
\(\sqrt{\operatorname{disc}(K)}\) and forces the subset to be full.
The scaling between complementary products has conjugate ratio
\(-1\).  Thus the two expansions with exponents
\(\pm\delta_{\mathfrak a}\) are exactly the two pullbacks indexed
by the complementary ramified products.

This combined function is nonconstant.  Indeed pass to a common
tame level and choose a Tate cusp at which all these prime kernels
are multiplicative.  They occur at distinct primes, disjoint from
the oriented split level, so the choices can be made simultaneously;
retain multiplicative type for those oriented cyclic groups too.
The identity translate has first positive Fourier coefficient
\(1\).  Every other ramified product has degree greater than one,
and its quotient expansion begins at a larger index.  No cancellation
of that coefficient is possible.
The injectivity of \(\Phi\) now contradicts the displayed zero sum.
Hence the conductor-one series is nonzero modulo \(p\) for each
orientation.

Finally apply Lemma~\ref{pair:trace}.  Each new factor
\eqref{eq:source-13} remains nonzero modulo \(p\): its true
\(\Gamma\)-exponent is nonzero by the tuple construction.
Its powers therefore vary nontrivially in \(\F_p[[t]]\), and the
displayed Laurent polynomial cannot vanish identically.
After setting the new variables to zero we obtain a product of
nonzero reductions.  The full series consequently has nonzero
reduction, proving \eqref{eq:source-14}.
\end{proof}

\subsection{Local comparison at high finite characters}

Fix a sufficiently high finite character \(\theta\) of \(\Gamma\)
and put
\[
 A_\theta=\Z_p[\theta][[\mathbf u]][1/p].
\]
First adjoin the unramified measure coefficients, obtaining
\[
 \widetilde A_\theta=\mathcal V^*[\theta][[\mathbf u]][1/p].
\]
Let \(A_\theta'\) be a finite coefficient extension of this ring
containing the Gauss coordinates \(\mu_{p^m}\) and the normal
resolvents for the fixed ramified local character fields used below.
More precisely, take the complete valuation ring of their
compositum with \(\Frac(\mathcal V^*[\theta])\), form its power-series
ring in \(\mathbf u\), and invert \(p\).
This finite extension is faithfully flat and may depend on
\(\theta\).  No bound on its ramification as \(\theta\) varies is
asserted or needed.
Since \(p\) is now inverted, these fixed nonzero constants, together
with the point-clearing factors, are units.  They will be allowed
only in this characteristic-zero comparison.  Divisibility of
series already in \(\widetilde A_\theta\) descends along this
faithfully flat extension, as does equality of their principal
ideals.  Proposition~\ref{pair:vertical-primitivity} supplies the
separate integral information at \(p\).

\begin{proposition}[Divisibility at high finite characters]
\label{pair:finite-divisibility}
For every sufficiently high finite \(\theta\),
\begin{equation}
 B_w(\theta,\mathbf u)\in
         L_w(\theta,\mathbf u)A_\theta'.
 \tag{15}\label{eq:source-15}
\end{equation}
With the same global character variable, the two nonzero quotients
for \(w\) and \(\bar w\) are associates in \(A_\theta'\).
\end{proposition}

To prove the proposition, we first identify the two \(p\)-local
Kummer lines and then express the quotient valuations through the
all-Kummer determinant.  The final step will be a sequence of
horizontal switches that preserves these valuations.

At a height-one prime of \(A_\theta\), the moving-place local
complexes are acyclic.  Their inertia detects their own variable
\(u_j\).  On the divisor \(u_j=0\), Frobenius-minus-one is still
generically invertible: its scalar contains the true exponent of
that moving prime, whose \(\theta\)-value has arbitrarily high
order.  For large \(\theta\), even setting the other variables to
zero cannot give either of the two elliptic eigenvalue exceptions.
The cardinality twist tends to \(1\).

At fixed \(q\in S^\circ\), replace the full condition by the
unramified condition after inverting \(p\).  The conditions are
exactly orthogonal, by the local comparison in
Section~\ref{sec:models}.  The product of the removed singular
determinants is \(D\) in \eqref{eq:source-11}, and is nonzero at
these tests.  The same replacement is available on horizontal
height-one localizations of \(R\).  Its strict determinant is
therefore \(L_w/D\), up to a unit.  This calculation uses the inverse
determinant convention throughout.

\begin{lemma}[The two local Kummer lines]\label{pair:local-lines}
At each \(p\)-adic place and after fixing \(\theta\) as above,
degree-one local cohomology splits over \(A_\theta\) into a free
Kummer line and its
free complementary line.  The paired Kummer conditions at conjugate
places are exactly orthogonal.  After extension to \(A_\theta'\),
the logarithm identifies the Kummer line with a free line, and its
value on the weighted CM point classes is
\eqref{eq:source-10}, with the factor \(c_E\) restored for
\(\log_{\omega_E}\).  These assertions hold for bad reduction at
\(p\) as well.
\end{lemma}

\begin{proof}
Let \(P_v/K_v\) be the finite totally ramified extension cut out by
the inertia part of \(\theta\), and separate its unramified part.
Write \(G_v=\Gal(P_v/K_v)\),
\(\mathcal O_\theta=\Z_p[\theta]\), and
\(\Delta=\Gal(D_\infty^{\rm loc}/K_v)\simeq\Z_p\), where
\(D_j^{\rm loc}/K_v\) is the unramified extension of degree \(p^j\).
Choose its arithmetic Frobenius generator \(\gamma_{\rm ur}\).
The Shapiro coefficient ring before finite-character projection is
\[
 \mathcal O_\theta[G_v][[\Delta]][1/p].
\]
Project by the character idempotent corresponding to the chosen
cohomological character.  The resulting scalar ring is
\(\Lambda_\theta=\mathcal O_\theta[[X]][1/p]\), with
\(1+X=\gamma_{\rm ur}\).
If \(e_{vj}\in\Z_p\) are the tame Frobenius exponents at \(v\), its
map to the comparison ring is
\[
 \Lambda_\theta\longrightarrow A_\theta,\qquad
 1+X\longmapsto
 \alpha_v\prod_{j=1}^d(1+u_j)^{e_{vj}},
 \quad \alpha_v=\theta_{\rm ur}(\Frob_v).
\]
The exponent vector is nonzero.  Divide it by its largest common
power \(p^a\), then complete the resulting primitive vector to a
\(\mathrm{GL}_d(\Z_p)\) coordinate change.  In the new group-power
coordinates the image is \(\alpha_v(1+y_1)^{p^a}\).
The equation
\[
 (1+y_1)^{p^a}=\alpha_v^{-1}(1+X)
\]
is monic and distinguished of degree \(p^a\) in \(y_1\).
It makes the one-variable substitution finite flat; adjoining the
remaining variables preserves flatness.  The marked local models
can therefore be base-changed through this map.

We first control lattices uniformly in \(j\), with constants allowed
to depend on \(\theta\).  A sufficiently small additive lattice
exponentiates because the ramification is fixed.  For an upper bound
pass to a fixed extension with semistable reduction.  The component
groups are bounded after adjoining \(P_v\) and making unramified
base changes.  Choose \(\theta\) nontrivial on sufficiently deep
inertia over this fixed field.  After multiplication by the bounded
component exponent, inertia acts trivially on the special points of
the identity component.  The Neron base-change map is surjective
on these geometric semiabelian points: its map on rational
auxiliary-adic Tate modules identifies the inertia invariants,
which are unchanged under finite extension once inertia is
unipotent.  Thus an inertia difference takes points into the formal
kernel.  A further fixed multiple takes them into a small logarithm
lattice.  On the projected eigenspace that difference is a nonzero
constant, giving the upper bound.  Use the unnormalized integral
projection to obtain the same bounded-exponent assertion for the
torsion logarithm kernel.  The argument applies to the inverse
character as well.  Trace-dual bounds follow from the bounded
different of the fixed ramified extension.

Integral Kummer and local duality on the Shapiro tower give an
exact sequence with ends the compact points and the dual of compact
points; transitions are norms and dual inclusions.  Compactness
makes the inverse limit exact.  One can identify both ends
explicitly from trace-compatible
integral normal generators in the unramified tower.  Choose their
successive traces compatibly and their total trace equal to \(1\).
Modulo \(p\), a unit augmentation generates the cyclic group-algebra
module, so these are normal generators integrally as well.
For the fixed ramified factor use a rational normal basis and its
bounded trace dual.  Before projection, the additive inverse limit
is the regular rank-one module over
\(\mathcal O_\theta[G_v][[\Delta]][1/p]\).  Apply the
finite-character idempotent to this additive module.  The projected
logarithm lattice bounds identify the projected compact point
module with it after inverting \(p\); the projected torsion kernel
is killed by the fixed bound already proved.  Thus the compact
point endpoint is free of rank one over \(\Lambda_\theta\).
The bounded trace-dual comparison gives the same conclusion for its
dual endpoint.  Applying the displayed flat substitution yields the
two free lines over \(A_\theta\).  The torsion bounds kill higher cohomology;
invariants vanish too.  The full local complex consequently splits.
Shapiro compatibility of the cup product, finite Kummer isotropy,
and local duality give exact orthogonality, including the
nullhomotopies in these split complexes.

Form the logarithmic group polynomials with the unnormalized
character weights before making the substitution.  The normal
resolvent on the ramified factor is nonzero; on the unramified
factors the compatible resolvent is a unit because its total trace
is \(1\).  Thus logarithm is the stated line isomorphism.
To identify transferred point classes, specialize first to bounded
cyclic quotients.  Freeness with coset bases identifies the
substitution with the corresponding finite unramified layer.
At sufficiently large stages the local ring-class fields contain
that layer, by the nonzero Frobenius exponent, and contain the
fixed ramified character field.  Trace the stage classes to these
fixed local fields and apply the finite-level logarithmic
calculation.  Membership in the Kummer line is detected on all
these finite quotients of its free complementary quotient.
The logarithmic expressions are detected in the same way.
The fixed-ramification bounds remain valid through the unramified
growth at \(p\), so passage to the bounded limits gives exactly
\eqref{eq:source-10}.
\end{proof}

\subsection{Horizontal switches and divisibility}

We now compare the order of the analytic series with the arithmetic
determinant at a height-one prime \(\mathfrak p\) of \(A_\theta\).
Write \(v_{\mathfrak p}\) for its normalized valuation.
Let \(Z\) be the Kummer class obtained by unnormalized weighted
corestriction of the CM points of conductor \(p^ms\).
One common fixed point multiple clears all descent requirements.
Use unramified conditions at the fixed non-\(p\) support, the two
Kummer lines at \(p\), and full conditions at moving primes.
Denote this complex by \(C_F\).  At \(\mathfrak p\) it is self-dual
under conjugate transport and the Weil pairing; the moving full
complexes are acyclic.

In a basis of the Kummer line at \(w\), let \(p_0\in A_\theta\)
be the localization of \(Z\).  By
Lemma~\ref{pair:local-lines} and \eqref{eq:source-10},
\[
 B_w(\theta,\mathbf u)
      =\text{unit in }A_\theta'\cdot p_0^2D.
\]
The constants included here are \(c_E\), the Gauss sum, and the
common clearing factors; they are units only because \(p\) has
already been inverted.  Both \(B_w(\theta,\mathbf u)\) and
\(L_w(\theta,\mathbf u)\) are nonzero for all sufficiently high
\(\theta\), by primitivity, \eqref{eq:source-1}, and bounded-series
preparation.  Hence \(p_0\ne0\).  The strict comparison shows that
the generic ranks of \(C_F\) are one in degrees one and two.
For its point and the calibrated class-functional dual, the
determinant triangle gives
\begin{equation}
 d\bigl(C_F,(Z,Z^\vee)\bigr)
     =2v_{\mathfrak p}(p_0)
           -v_{\mathfrak p}(L_w/D).
 \tag{16}\label{eq:source-16}
\end{equation}
This is the line comparison of
\cite[Section~3, equation (D)]{B}, with the integral
pairings and inverse determinant of Section~\ref{sec:models}.

\begin{lemma}[Inert derivative coordinates]\label{pair:derivatives}
There are fresh inert good prime sequences \(\ell\), with
\(\ell\to-1\) and Tate Frobenius tending to
\(J=\rho_T(\tau)\) for complex conjugation \(\tau\), for which
the limiting local degree-one cohomology splits into finite and
transverse planes.  Their mixed pairing is, up to a unit,
\(\langle x,Jy\rangle\), and both planes are self-annihilating.
For each fixed finite sequence of switches, there is a compatible
family of descended weighted Kummer classes \(Z_I\), indexed by
the subsets occurring in that sequence, with \(Z_\varnothing=Z\).
The class \(Z_I\) uses the product of the ring-class derivatives
\(D_\ell\) for \(\ell\in I\); its local conditions are transverse
at \(I\) and the original conditions at the other places.
These classes satisfy
\[
 f_{\rm loc}(Z_{I\cup\ell})=0,\qquad
 s_{\rm loc}(Z_{I\cup\ell})=Jf_{\rm loc}(Z_I).
\]
All prior local conditions can be retained, with one common point
scaling for the classes in the fixed finite sequence.
\end{lemma}

\begin{proof}
Choose the Frobenius sequence \(\gamma_i\to\tau\) on the full Tate
module.  Then \(a_f(\ell)\to0\), while
\(\gamma_i^2\) acts trivially on the ring-class fields by the
dihedral relation.  The old characters are unramified at \(\ell\).
In the cyclic inertia resolution the residue conjugating power
tends to \(1\).  Thus the limiting full local cohomology is split
free.  Evaluation on \(\gamma_i^2\) and on a tame generator gives
the finite and transverse coordinates.  Retain invariants in the
lower condition and omit \(H^2\) in the upper condition, as required
by the square-switch lemma.  Conjugate transport through
\(\gamma_i\), together with local reciprocity, gives mixed pairing
\(\langle x,Jy\rangle\).  Pure scalar cups vanish and the mixed
scalar cup is a unit.  These are exactly the Weil-plane hypotheses
of that lemma.

Here is the derivative calculation, following
\cite[Section~9, \emph{Inert derivative identity}, equation (K)]{B}, with the
odd-\(p\) precision specified.  The relative ring-class order at
an inert prime is \(\ell+1\), since the only order units are signs.
For its generator \(\sigma_\ell\) use
\[
 D_\ell=\sum_{j=1}^{\ell}j\sigma_\ell^j.
\]
Writing \(\operatorname{Tr}_\ell=\sum_{j=0}^{\ell}\sigma_\ell^j\),
the group-ring identity is
\[
 (\sigma_\ell-1)D_\ell=(\ell+1)-\operatorname{Tr}_\ell.
\]
If \(Y\) is the upper point and \(X\) the lower point, the Hecke
trace gives \(a_\ell X\), after the common cusp multiple, and
supersingular reduction gives
\(\widetilde Y=F_\ell\widetilde X\).
Choose growing precision \(p^a\) slowly enough that
\(p^a\mid \ell+1\), \(p^a\mid a_\ell\), and
\(\gamma_i^2=1\) modulo \(p^a\).
The group-ring derivative identity then makes the upper Kummer
class invariant to that precision.

Inflation and restriction descend it after a fixed multiple.
The needed uniform bound on invariant torsion follows from
Bogomolov's homothety theorem \cite{Bogomolov}: take
\(h\in G_K\) scalar sufficiently close to \(1\), with square
different from \(1\).  The element \(h\tau h\tau\) kills every
ring-class action and acts on the Tate module by that scalar
squared.  It bounds torsion over all fields in use.  At the old
\(p\)-places, descent is unramified and retains the Kummer
condition after a fixed component multiple.  Lang's theorem handles
the connected special points; the formal-group assertion follows
from induced additive successive quotients and passage to the
limit.  Component groups remain bounded under unramified base
change above the fixed ramified extensions.  The same argument
retains unramified conditions at the old non-\(p\) support.  Include
all these multiples in one scale for every switched class.

To verify the coordinates, suppress that scale and choose
\(p^aQ=X\), \(p^aW=D_\ell Y\).  If \(b\) is the descended upper
cocycle, the crossed identity makes
\((\sigma_\ell-1)W-b(\sigma_\ell)\) fixed by the upper-field
subgroup.  Multiplication by \(p^a\) gives
\((\ell+1)Y-a_\ell X\), so it differs from
\[
 ((\ell+1)/p^a)Y-(a_\ell/p^a)X
\]
by bounded-exponent torsion.  After reduction, inertia is trivial.
Thus, to the retained precision, \(b(\sigma_\ell)\) gives
\[
 (a_\ell-(\ell+1)F_\ell)\widetilde Q
       =(F_\ell-a_\ell)(F_\ell^2-1)\widetilde Q,
\]
where the equality is Eichler--Shimura.  The lower Kummer cocycle,
evaluated on the residue Frobenius \(\gamma_i^2\), has finite
coordinate
\[
 f_{\rm loc}(Z_I)=(F_\ell^2-1)\widetilde Q
\]
at the retained precision.  The displayed value of
\(b(\sigma_\ell)\) is the transverse coordinate of the upper class.
Consequently it equals
\((F_\ell-a_\ell)f_{\rm loc}(Z_I)\).
Passing to the prescribed limit gives \(F_\ell-a_\ell\to J\),
which proves the asserted transverse identity.
For the upper finite coordinate,
\(\widetilde{D_\ell Y}\) has an \(F_\ell^2\)-fixed
\(p^a\)-division obtained using
\(\ell(\ell+1)/(2p^a)\) times \(F_\ell\widetilde X\);
the ambiguity is killed by \(F_\ell^2-1\).  Reduction detects
prime-to-\(\ell\) torsion, proving the zero finite coordinate.
The calculation applies on each translate before weighted
corestriction, with any earlier derivative factors already present.
It therefore proves the two identities simultaneously at all
previous switches and with the identical common scale.
\end{proof}

\begin{proof}[Proof of Proposition~\ref{pair:finite-divisibility}]
We show that \eqref{eq:source-16} is nonnegative at every height-one
test.  If fiber degree one has dimension greater than one, choose
two independent classes, one being a primitive reduction from the
generic line.  They admit a rank-two evaluation on the finite
plane at a suitable inert prime.

We spell out this detection because the same assertion must hold
after previous switches.  Lemma~\ref{models:evaluations}
realizes the fiber classes on the product
of the stage Galois groups.  Restriction to the joint kernel of
the full Tate and ring-class actions detects them: on the image
quotient the central element \(h\tau h\tau\) just used has scalar
difference from \(1\) invertible.  Global invariants vanish.
The Selmer classes inject into unrestricted cohomology by the
local hypotheses, and the complexes have amplitude \(1,2\) by
duality.  Let \(\kappa(\mathfrak p)\) be the testing residue field,
and write
\[
 W=(V\otimes\kappa(\mathfrak p))\chi,\qquad
 W^c\simeq(V\otimes\kappa(\mathfrak p))\chi^{-1}.
\]
Conjugate transport makes these the two summands of the induced
representation after adjoining diagonal \(\tau\); the element
\(\tau\) exchanges them through the elliptic involution \(J\).
They are absolutely simple and nonisomorphic: their scalar
twists have distinct determinants because the inertia character
has sufficiently high order.  For the two independent classes
and their conjugates, joint evaluations on the kernel therefore
span \(W^2\oplus(W^c)^2\).
For an element \(\eta_b\) of that kernel, evaluation on
\((\tau\eta_b)^2\) is additive in the kernel evaluations.
On a tuple \((a_1,a_2;b_1,b_2)\), the operator \(1+\tau\)
followed by projection to \(W^2\) gives
\[
 (a_1+Jb_1,\ a_2+Jb_2).
\]
This map is surjective, so its evaluations still span both copies
needed for a rank-two test.
An additive subgroup contains the integer grid generated by any
finite spanning set; in characteristic zero this grid cannot lie
in a proper rank-defect locus.  A rank-two choice follows.
Finite Chebotarev realizes these evaluations to increasing
precision by primes whose Frobenius agrees with
\(\tau\eta_b\).

Apply Lemma~\ref{models:square-switch}, using
Lemma~\ref{pair:derivatives}.  The fiber dimension drops by two;
the generic ranks and the valuation of the point/class-functional
tensor are preserved.  The fiber dimension cannot become zero:
upper semicontinuity bounds it below by the generic rank one.
Thus the repeated two-dimensional drops end at dimension one.
There the complex is homotopy equivalent to two free lines with
zero differential and perfect duality.  Its integral point tensor
has nonnegative determinant valuation.  All fixed scales are
units at this test.  Thus \eqref{eq:source-16} is nonnegative.
The height-one criterion in the UFD \(A_\theta\), and its
coefficient extension, gives
\[
 p_0^2/(L_w/D)\in A_\theta',
\]
which is \eqref{eq:source-15}.

For the association assertion, keep the global character unchanged
and perform the same comparison at \(\bar w\).
At conductor \(p^ms\), for large \(m\), the canonical connected
quotient lowers the Serre--Tate parameter order by \(p\).
It is the unique ascending edge in the CM isogeny graph.
The required fixed number of cyclic \(p\)-level edges therefore
gives the same level along the class orbit for the two choices
of \(p\)-adic place.  Their points agree up to the non-\(p\)
Atkin--Lehner operations and translations already considered.
The weighted classes differ only by units in \(A_\theta\);
the cuspidal differences are killed.  Hence the all-Kummer
complex and its tensor in \eqref{eq:source-16} are the same up to
units at every height-one test.  The two quotient valuations
coincide, and the quotients are associates.  Contracting the equal
principal ideals along the finite faithfully flat coefficient
extension gives the same association over
\(\widetilde A_\theta\).
The high-conductor
condition is essential here; this argument does not compare two
arbitrarily normalized logarithmic lines at a low-conductor
specialization.
\end{proof}

\begin{corollary}[The integral comparison quotient]
\label{pair:integral-quotient}
Using the descended analytic ideals, the fraction
\[
 U_w=B_w/L_w
\]
is represented by an integral series, and \(p\nmid U_w\).
Removing any newly added tame variables specializes it to the
old quotient times an integral unit.  This includes removal of
all tame variables, leaving the genuine anticyclotomic line.
\end{corollary}

\begin{proof}
By \eqref{eq:source-1}, \(L_w\) is nonzero modulo \((p,t)\).
Lemma~\ref{models:series-division}
therefore applies to Proposition~\ref{pair:finite-divisibility}.
For clarity, make \(L_w\) distinguished in one tame variable by
an integral group-power change and divide \(B_w\) by it
integrally.  At every sufficiently high character of \(t\), the
remainder is zero after inverting \(p\) and extending coefficients
to \(A_\theta'\).  Faithful flatness descends this vanishing to
\(\widetilde A_\theta\); the monic quotient has
no \(p\)-torsion, so it is zero integrally there.  Each bounded
coefficient of the remainder vanishes at all those characters
and hence vanishes identically by one-variable Weierstrass
preparation.  This proves integral divisibility.  Since
\(p\nmid B_w\), the integral quotient has no factor \(p\).

Lemma~\ref{pair:trace} restores the same squared factor in
numerator and denominator on setting each new variable to zero.
That factor is nonzero, and the determinant comparison unit is
integral.  Cancelling in the fraction field proves the asserted
specialization identity.  Integrality of the larger quotient
then makes the specialization integral as well.  There is no
loss of precision in this step.
\end{proof}

We have obtained the integral quotient, its exact specialization
law, and association of its high-character specializations at the
two \(p\)-adic places.  Its possible horizontal divisors remain.
The next two sections construct the theta implication needed to
exclude them.
